\part{Engine architecture}

\chapter{Shared contracts and why the repo is contract-first}
\label{ch:shared-contracts}
The repository is organised around shared contracts rather than independent language ports. The specs define behaviour; artifacts carry data; implementations consume the same contract.

\begin{figure}[htbp]
\centering
\resizebox{0.94\textwidth}{!}{%
\begin{tikzpicture}[node distance=8mm,>=Latex]
\tikzstyle{box}=[rounded rectangle,draw=EngineDark!70,fill=EngineGrey,minimum width=31mm,minimum height=9mm,align=center,font=\small]
\tikzstyle{blue}=[rounded rectangle,draw=EngineBlue!80,fill=EngineBlue!7,minimum width=31mm,minimum height=9mm,align=center,font=\small]
\node[blue] (specs) {shared specs};
\node[box,below=of specs] (tools) {conversion tools};
\node[box,below=of tools] (artifacts) {converted artifacts};
\node[box,below left=12mm and 14mm of artifacts] (pt) {PyTorch reference};
\node[box,below=12mm of artifacts] (cpp) {C++ engine};
\node[box,below right=12mm and 14mm of artifacts] (rs) {Rust engine};
\node[blue,below=18mm of cpp] (val) {validation and benchmarks};
\draw[->,thick] (specs)--(tools);
\draw[->,thick] (tools)--(artifacts);
\draw[->,thick] (artifacts)--(pt);
\draw[->,thick] (artifacts)--(cpp);
\draw[->,thick] (artifacts)--(rs);
\draw[->,thick,EngineGreen!70!black] (pt)--(val);
\draw[->,thick,EngineGreen!70!black] (cpp)--(val);
\draw[->,thick,EngineGreen!70!black] (rs)--(val);
\end{tikzpicture}
}
\caption{Contract-first architecture: three runtimes consume one artifact and one validation contract.}
\end{figure}


\begin{intuitionbox}{Intuition}
A shared contract is what makes cross-language comparison meaningful.
\end{intuitionbox}

\begin{definitionbox}{Mathematical or contract statement}
The shared contract includes model config, tensor names, artifact layout, cache semantics, test-vector schema, CLI output, and benchmark separation.
\end{definitionbox}

\begin{implementationbox}{Implementation interpretation}
Put artifact validation before model execution. A runtime that cannot prove the artifact is valid should not produce logits.
\end{implementationbox}

\begin{verificationbox}{How to verify this works}
\begin{itemize}
\item Check shapes and dtypes at the boundary.
\item Compare deterministic outputs against PyTorch golden vectors.
\item Record tolerance, seed, model artifact, and backend.
\item Do not benchmark until validation passes.
\end{itemize}
\end{verificationbox}

\begin{warningbox}{Common failure modes}
\begin{itemize}
\item Letting framework defaults choose dtype or layout.
\item Testing only a batch size of one.
\item Depending on upstream checkpoint names in runtime code.
\item Treating benchmark output as validation.
\end{itemize}
\end{warningbox}

\begin{chapterexercises}
\item State the central invariant in one sentence.\\
\textit{Hint.} Use the conceptual guide vocabulary.
\item Construct a toy example with small shapes.\\
\textit{Hint.} Use $B=2$, $T=3$, $H=8$ where possible.
\item Name a validation test that would catch a plausible bug.\\
\textit{Hint.} Prefer generated golden vectors to handwritten long arrays.
\end{chapterexercises}

\chapter{Model configuration}
\label{ch:model-config}
The model config is internal to the engine. It is not a Hugging Face config.


\begin{table}[htbp]\centering\small
\begin{tabularx}{\textwidth}{L{0.30\textwidth}Y}
\toprule
Field group & Required semantics \\
\midrule
Identity & \code{model_family}, \code{model_name}, \code{architecture}. \\
Sizes & \code{vocab_size}, \code{max_position_embeddings}, \code{hidden_size}, layers, heads, KV heads, intermediate size. \\
Module choices & Activation, norm type, positional encoding, MLP type, attention type. \\
Flags & Attention bias, MLP bias, tied embeddings. \\
Tokens and dtype & BOS, EOS, and phase-0 \dtype{float32}. \\
Nullable later fields & \code{rope_theta}, \code{rope_scaling}, initializer range. \\
\bottomrule
\end{tabularx}
\caption{Model configuration groups.}
\end{table}


\begin{intuitionbox}{Intuition}
Configuration is the type signature of the model.
\end{intuitionbox}

\begin{definitionbox}{Mathematical or contract statement}
Phase-0 GPT-2 requires \code{architecture=decoder_only}, \dtype{float32}, learned absolute positions, MHA, and \code{num_key_value_heads == num_attention_heads}.
\end{definitionbox}

\begin{implementationbox}{Implementation interpretation}
Parse into a strongly checked \code{ModelConfig}. Compute derived values such as \code{head_dim} only after validation.
\end{implementationbox}

\begin{verificationbox}{How to verify this works}
\begin{itemize}
\item Check shapes and dtypes at the boundary.
\item Compare deterministic outputs against PyTorch golden vectors.
\item Record tolerance, seed, model artifact, and backend.
\item Do not benchmark until validation passes.
\end{itemize}
\end{verificationbox}

\begin{warningbox}{Common failure modes}
\begin{itemize}
\item Letting framework defaults choose dtype or layout.
\item Testing only a batch size of one.
\item Depending on upstream checkpoint names in runtime code.
\item Treating benchmark output as validation.
\end{itemize}
\end{warningbox}

\begin{chapterexercises}
\item State the central invariant in one sentence.\\
\textit{Hint.} Use the conceptual guide vocabulary.
\item Construct a toy example with small shapes.\\
\textit{Hint.} Use $B=2$, $T=3$, $H=8$ where possible.
\item Name a validation test that would catch a plausible bug.\\
\textit{Hint.} Prefer generated golden vectors to handwritten long arrays.
\end{chapterexercises}

\chapter{Tensor naming and weight artifact design}
\label{ch:tensor-naming}
Converted artifacts use canonical internal tensor names. Runtime loaders must not depend on upstream checkpoint names.


{\scriptsize
\begin{longtable}{L{0.42\textwidth}L{0.28\textwidth}L{0.20\textwidth}}
\caption{Required GPT-2 internal tensor names.}\\
\toprule
Tensor & Shape & Role \\
\midrule
\endfirsthead
\toprule
Tensor & Shape & Role \\
\midrule
\endhead
\tensorname{tok_embeddings.weight} & $[V,H]$ & Token table \\
\tensorname{pos_embeddings.weight} & $[P,H]$ & Positions \\
\tensorname{blocks.i.ln_1.weight} and \tensorname{bias} & $[H]$ & First norm \\
\tensorname{blocks.i.attn.qkv.weight} & $[3H,H]$ & Fused QKV \\
\tensorname{blocks.i.attn.qkv.bias} & $[3H]$ & QKV bias \\
\tensorname{blocks.i.attn.out_proj.weight} & $[H,H]$ & Output projection \\
\tensorname{blocks.i.attn.out_proj.bias} & $[H]$ & Output bias \\
\tensorname{blocks.i.ln_2.weight} and \tensorname{bias} & $[H]$ & Second norm \\
\tensorname{blocks.i.mlp.fc_in.weight} & $[F,H]$ & MLP expansion \\
\tensorname{blocks.i.mlp.fc_in.bias} & $[F]$ & MLP bias \\
\tensorname{blocks.i.mlp.fc_out.weight} & $[H,F]$ & MLP contraction \\
\tensorname{blocks.i.mlp.fc_out.bias} & $[H]$ & MLP bias \\
\tensorname{ln_f.weight} and \tensorname{bias} & $[H]$ & Final norm \\
\tensorname{lm_head.weight} & $[V,H]$ & Explicit LM head \\
\bottomrule
\end{longtable}}


\begin{intuitionbox}{Intuition}
Names are part of the engine ABI.
\end{intuitionbox}

\begin{definitionbox}{Mathematical or contract statement}
GPT-2 uses top-level prefixes \tensorname{tok_embeddings}, \tensorname{pos_embeddings}, \tensorname{blocks}, \tensorname{ln_f}, and \tensorname{lm_head}.
\end{definitionbox}

\begin{implementationbox}{Implementation interpretation}
The converter owns upstream-name knowledge. Runtime loaders use only internal names and required shapes.
\end{implementationbox}

\begin{verificationbox}{How to verify this works}
\begin{itemize}
\item Check shapes and dtypes at the boundary.
\item Compare deterministic outputs against PyTorch golden vectors.
\item Record tolerance, seed, model artifact, and backend.
\item Do not benchmark until validation passes.
\end{itemize}
\end{verificationbox}

\begin{warningbox}{Common failure modes}
\begin{itemize}
\item Letting framework defaults choose dtype or layout.
\item Testing only a batch size of one.
\item Depending on upstream checkpoint names in runtime code.
\item Treating benchmark output as validation.
\end{itemize}
\end{warningbox}

\begin{chapterexercises}
\item State the central invariant in one sentence.\\
\textit{Hint.} Use the conceptual guide vocabulary.
\item Construct a toy example with small shapes.\\
\textit{Hint.} Use $B=2$, $T=3$, $H=8$ where possible.
\item Name a validation test that would catch a plausible bug.\\
\textit{Hint.} Prefer generated golden vectors to handwritten long arrays.
\end{chapterexercises}

\chapter{SafeTensors and conversion into internal artifacts}
\label{ch:safetensors-conversion}
The project supports direct SafeTensors loading for inspection and conversion from SafeTensors or upstream checkpoints into internal artifacts. \citep{safetensorsDocs}


\begin{figure}[htbp]
\centering
\resizebox{0.94\textwidth}{!}{%
\begin{tikzpicture}[node distance=7mm,>=Latex]
\tikzstyle{filebox}=[rounded rectangle,draw=EnginePurple!80,fill=EnginePurple!6,minimum width=45mm,minimum height=9mm,align=center,font=\small]
\node[filebox] (cfg) {config.json};
\node[filebox,below=of cfg] (meta) {metadata.json};
\node[filebox,below=of meta] (idx) {weights.index.json};
\node[filebox,below=of idx] (bin) {weights.bin};
\draw[->,thick] (cfg)--(meta);
\draw[->,thick] (meta)--(idx);
\draw[->,thick] (idx)--(bin);
\node[rounded rectangle,draw=EngineOrange!80,fill=EngineOrange!7,text width=55mm,align=left,right=16mm of idx,font=\small] {Runtime code trusts internal tensor names and authoritative byte offsets, not upstream checkpoint names.};
\end{tikzpicture}
}
\caption{Internal artifact format and its authority boundaries.}
\end{figure}

\begin{lstlisting}[caption={Schematic weights.index.json entry.}]
{
  "blocks.0.attn.qkv.weight": {
    "dtype": "float32",
    "shape": [2304, 768],
    "offset": 1048576,
    "nbytes": 7077888,
    "source_name": "transformer.h.0.attn.c_attn.weight"
  }
}
\end{lstlisting}


\begin{intuitionbox}{Intuition}
Direct loading is convenient; conversion is the contract boundary.
\end{intuitionbox}

\begin{definitionbox}{Mathematical or contract statement}
A valid artifact has \file{config.json}, \file{metadata.json}, \file{weights.index.json}, and \file{weights.bin}. Index offsets and byte counts are authoritative.
\end{definitionbox}

\begin{implementationbox}{Implementation interpretation}
Implement \code{SafeTensorLoader} for direct reads, \code{ArtifactConverter} for mapping into internal names, and \code{WeightLoader} for runtime artifacts.
\end{implementationbox}

\begin{verificationbox}{How to verify this works}
\begin{itemize}
\item Check shapes and dtypes at the boundary.
\item Compare deterministic outputs against PyTorch golden vectors.
\item Record tolerance, seed, model artifact, and backend.
\item Do not benchmark until validation passes.
\end{itemize}
\end{verificationbox}

\begin{warningbox}{Common failure modes}
\begin{itemize}
\item Letting framework defaults choose dtype or layout.
\item Testing only a batch size of one.
\item Depending on upstream checkpoint names in runtime code.
\item Treating benchmark output as validation.
\end{itemize}
\end{warningbox}

\begin{chapterexercises}
\item State the central invariant in one sentence.\\
\textit{Hint.} Use the conceptual guide vocabulary.
\item Construct a toy example with small shapes.\\
\textit{Hint.} Use $B=2$, $T=3$, $H=8$ where possible.
\item Name a validation test that would catch a plausible bug.\\
\textit{Hint.} Prefer generated golden vectors to handwritten long arrays.
\end{chapterexercises}

\chapter{Tokenizers and embedders}
\label{ch:tokenizers-embedders}
Tokenizers convert text to ids and ids back to text. Tokenizer validation is separate from model-math validation.



\begin{intuitionbox}{Intuition}
The model consumes token ids, not strings.
\end{intuitionbox}

\begin{definitionbox}{Mathematical or contract statement}
Token embeddings map $I\in\N^{B\times T}$ to $X\in\R^{B\times T\times H}$ by table lookup; learned positions add a row per position.
\end{definitionbox}

\begin{implementationbox}{Implementation interpretation}
For correctness vectors, \code{input_ids} are authoritative and prompt text is optional metadata.
\end{implementationbox}

\begin{verificationbox}{How to verify this works}
\begin{itemize}
\item Check shapes and dtypes at the boundary.
\item Compare deterministic outputs against PyTorch golden vectors.
\item Record tolerance, seed, model artifact, and backend.
\item Do not benchmark until validation passes.
\end{itemize}
\end{verificationbox}

\begin{warningbox}{Common failure modes}
\begin{itemize}
\item Letting framework defaults choose dtype or layout.
\item Testing only a batch size of one.
\item Depending on upstream checkpoint names in runtime code.
\item Treating benchmark output as validation.
\end{itemize}
\end{warningbox}

\begin{chapterexercises}
\item State the central invariant in one sentence.\\
\textit{Hint.} Use the conceptual guide vocabulary.
\item Construct a toy example with small shapes.\\
\textit{Hint.} Use $B=2$, $T=3$, $H=8$ where possible.
\item Name a validation test that would catch a plausible bug.\\
\textit{Hint.} Prefer generated golden vectors to handwritten long arrays.
\end{chapterexercises}

\chapter{Decoder-only model architecture}
\label{ch:decoder-only}
The engine targets decoder-only autoregressive models. GPT-2 proves the baseline dense decoder stack.



\begin{intuitionbox}{Intuition}
A decoder sees past tokens and never future tokens.
\end{intuitionbox}

\begin{definitionbox}{Mathematical or contract statement}
The model maps ids to logits $Z\in\R^{B\times T\times V}$ through embeddings, blocks, final normalisation, and an LM head.
\end{definitionbox}

\begin{implementationbox}{Implementation interpretation}
Keep block variants behind configuration and model-family adapters.
\end{implementationbox}

\begin{verificationbox}{How to verify this works}
\begin{itemize}
\item Check shapes and dtypes at the boundary.
\item Compare deterministic outputs against PyTorch golden vectors.
\item Record tolerance, seed, model artifact, and backend.
\item Do not benchmark until validation passes.
\end{itemize}
\end{verificationbox}

\begin{warningbox}{Common failure modes}
\begin{itemize}
\item Letting framework defaults choose dtype or layout.
\item Testing only a batch size of one.
\item Depending on upstream checkpoint names in runtime code.
\item Treating benchmark output as validation.
\end{itemize}
\end{warningbox}

\begin{chapterexercises}
\item State the central invariant in one sentence.\\
\textit{Hint.} Use the conceptual guide vocabulary.
\item Construct a toy example with small shapes.\\
\textit{Hint.} Use $B=2$, $T=3$, $H=8$ where possible.
\item Name a validation test that would catch a plausible bug.\\
\textit{Hint.} Prefer generated golden vectors to handwritten long arrays.
\end{chapterexercises}

\chapter{Prefill, decode, and generation loops}
\label{ch:prefill-decode}
Prefill processes a prompt and builds cache state. Decode processes incremental tokens using cached state.

\begin{figure}[htbp]
\centering
\resizebox{0.94\textwidth}{!}{%
\begin{tikzpicture}[node distance=8mm,>=Latex]
\tikzstyle{proc}=[rounded rectangle,draw=EngineBlue!80,fill=EngineBlue!6,minimum width=31mm,minimum height=9mm,align=center,font=\small]
\tikzstyle{state}=[rounded rectangle,draw=EnginePurple!80,fill=EnginePurple!6,minimum width=31mm,minimum height=9mm,align=center,font=\small]
\node[proc] (prompt) {prompt ids};
\node[proc,right=of prompt] (prefill) {prefill};
\node[state,right=of prefill] (cache) {KV cache};
\node[proc,below=of cache] (decode) {decode};
\node[proc,left=of decode] (sample) {sample id};
\node[state,right=of decode] (cache2) {cache append};
\draw[->,thick] (prompt)--(prefill);
\draw[->,thick] (prefill)--(cache);
\draw[->,thick] (cache)--(decode);
\draw[->,thick] (decode)--(sample);
\draw[->,thick] (decode)--(cache2);
\end{tikzpicture}
}
\caption{Prefill builds cache state; decode reuses and appends it.}
\end{figure}


\begin{intuitionbox}{Intuition}
Generation is stateful and must be tested as stateful.
\end{intuitionbox}

\begin{definitionbox}{Mathematical or contract statement}
Cache equivalence states that full recomputation on prompt plus token equals prefill on the prompt followed by cached decode within tolerance.
\end{definitionbox}

\begin{implementationbox}{Implementation interpretation}
Expose \code{Prefill}, \code{Decode}, and \code{GenerationLoop}. Keep stop conditions and sampler policy explicit.
\end{implementationbox}

\begin{verificationbox}{How to verify this works}
\begin{itemize}
\item Check shapes and dtypes at the boundary.
\item Compare deterministic outputs against PyTorch golden vectors.
\item Record tolerance, seed, model artifact, and backend.
\item Do not benchmark until validation passes.
\end{itemize}
\end{verificationbox}

\begin{warningbox}{Common failure modes}
\begin{itemize}
\item Letting framework defaults choose dtype or layout.
\item Testing only a batch size of one.
\item Depending on upstream checkpoint names in runtime code.
\item Treating benchmark output as validation.
\end{itemize}
\end{warningbox}

\begin{chapterexercises}
\item State the central invariant in one sentence.\\
\textit{Hint.} Use the conceptual guide vocabulary.
\item Construct a toy example with small shapes.\\
\textit{Hint.} Use $B=2$, $T=3$, $H=8$ where possible.
\item Name a validation test that would catch a plausible bug.\\
\textit{Hint.} Prefer generated golden vectors to handwritten long arrays.
\end{chapterexercises}

\chapter{KV-cache abstraction and cache policies}
\label{ch:kv-cache}
The KV cache stores recurrent decode state per layer. Phase 0 uses append-only cache; later phases add sliding-window policies.

\begin{figure}[htbp]
\centering
\resizebox{0.94\textwidth}{!}{%
\begin{tikzpicture}[>=Latex]
\draw[fill=EngineBlue!8,draw=EngineBlue!70] (0,0) rectangle (7.2,2.8);
\draw[fill=EngineTeal!10,draw=EngineTeal!70] (0.5,0.45) rectangle (7.7,3.25);
\draw[fill=EngineGreen!10,draw=EngineGreen!70] (1.0,0.9) rectangle (8.2,3.7);
\node at (4.2,2.35) {logical KV cache};
\draw[<->,thick] (1.0,0.75)--(8.2,0.75) node[midway,below] {seq};
\draw[<->,thick] (8.45,0.9)--(8.45,3.7) node[midway,right] {head dim};
\node[font=\small] at (4.1,-0.4) {$[B,H_{kv},S,d_h]$ logical layout; physical layout may differ.};
\end{tikzpicture}
}
\caption{Canonical KV-cache logical layout.}
\end{figure}


\begin{intuitionbox}{Intuition}
A cache is a visibility rule plus stored tensors.
\end{intuitionbox}

\begin{definitionbox}{Mathematical or contract statement}
For each layer \(\ell\), keys and values have logical shape \(K_\ell,V_\ell\in\R^{B\times H_{kv}\times S\times d_h}\). In phase 0, the visible span for decode step \(t\) is the full accumulated prefix \(s\leq S+t\); sliding-window policies later replace this with a bounded visibility predicate.
\end{definitionbox}

\begin{implementationbox}{Implementation interpretation}
Implement \code{KVCache}, \code{CachePolicy}, \code{AppendOnlyCache}, and \code{SlidingWindowCache}; only append-only is active in phase 0.
\end{implementationbox}

\begin{verificationbox}{How to verify this works}
\begin{itemize}
\item Check shapes and dtypes at the boundary.
\item Compare deterministic outputs against PyTorch golden vectors.
\item Record tolerance, seed, model artifact, and backend.
\item Do not benchmark until validation passes.
\end{itemize}
\end{verificationbox}

\begin{warningbox}{Common failure modes}
\begin{itemize}
\item Letting framework defaults choose dtype or layout.
\item Testing only a batch size of one.
\item Depending on upstream checkpoint names in runtime code.
\item Treating benchmark output as validation.
\end{itemize}
\end{warningbox}

\begin{chapterexercises}
\item State the central invariant in one sentence.\\
\textit{Hint.} Use the conceptual guide vocabulary.
\item Construct a toy example with small shapes.\\
\textit{Hint.} Use $B=2$, $T=3$, $H=8$ where possible.
\item Name a validation test that would catch a plausible bug.\\
\textit{Hint.} Prefer generated golden vectors to handwritten long arrays.
\end{chapterexercises}

\chapter{CLI, validation, and benchmarking contracts}
\label{ch:cli-contracts}
The CLI contract makes all runtimes comparable from automation. Required commands are run, generate, test, and bench.


\begin{lstlisting}[caption={Common CLI operations.}]
engine run --model <model_dir> --input-ids 15496,995
engine generate --model <model_dir> --input-ids 15496,995 --max-new-tokens 16
engine test --model <model_dir> --test-vector <path>
engine bench --model <model_dir> --suite <suite_config>
\end{lstlisting}


\begin{intuitionbox}{Intuition}
A CLI is a reproducibility surface.
\end{intuitionbox}

\begin{definitionbox}{Mathematical or contract statement}
Automation-facing commands must support JSON output and nonzero exit codes for failure. Correctness workflows should prefer input ids.
\end{definitionbox}

\begin{implementationbox}{Implementation interpretation}
The model path points to a converted artifact directory, never an upstream checkpoint cache.
\end{implementationbox}

\begin{verificationbox}{How to verify this works}
\begin{itemize}
\item Check shapes and dtypes at the boundary.
\item Compare deterministic outputs against PyTorch golden vectors.
\item Record tolerance, seed, model artifact, and backend.
\item Do not benchmark until validation passes.
\end{itemize}
\end{verificationbox}

\begin{warningbox}{Common failure modes}
\begin{itemize}
\item Letting framework defaults choose dtype or layout.
\item Testing only a batch size of one.
\item Depending on upstream checkpoint names in runtime code.
\item Treating benchmark output as validation.
\end{itemize}
\end{warningbox}

\begin{chapterexercises}
\item State the central invariant in one sentence.\\
\textit{Hint.} Use the conceptual guide vocabulary.
\item Construct a toy example with small shapes.\\
\textit{Hint.} Use $B=2$, $T=3$, $H=8$ where possible.
\item Name a validation test that would catch a plausible bug.\\
\textit{Hint.} Prefer generated golden vectors to handwritten long arrays.
\end{chapterexercises}

\chapter{High-level PyTorch/Rust/C++ interface design}
\label{ch:interfaces}
The same conceptual interfaces appear in PyTorch, C++, and Rust even when language idioms differ.


{\scriptsize
\begin{longtable}{L{0.42\textwidth}L{0.52\textwidth}}
\caption{Abstract components required by the conceptual guide.}\\
\toprule
Component & Responsibility \\
\midrule
\endfirsthead
\toprule
Component & Responsibility \\
\midrule
\endhead
\texttt{ModelConfig} & Validated engine-native model configuration. \\
\texttt{WeightLoader}, \texttt{SafeTensorLoader}, \texttt{ArtifactConverter} & Direct SafeTensors reading where useful and conversion into internal artifacts. \\
\texttt{Tensor}, \texttt{Device}, \texttt{DType} & Shape, dtype, layout, storage, and placement. \\
\texttt{Tokenizer}, \texttt{Embedder}, \texttt{TokenEmbedding} & Text boundary and embedding lookup. \\
\texttt{PositionalEncoding}, \texttt{LearnedAbsolutePositionEncoding}, \texttt{RoPE} & Position strategies. \\
\texttt{Attention}, \texttt{SelfAttention}, \texttt{MultiHeadAttention}, \texttt{MQA}, \texttt{GQA}, \texttt{AttentionKernel} & Attention semantics and backend kernels. \\
\texttt{MLPBlock}, \texttt{GELUFFN}, \texttt{SwiGLU}, \texttt{GeGLU}, \texttt{Activation} & Feed-forward families and activations. \\
\texttt{LayerNorm}, \texttt{RMSNorm}, \texttt{ResidualBlock}, \texttt{TransformerBlock} & Normalisation and block composition. \\
\texttt{MambaBlock} & Hybrid or recurrent block abstraction. \\
\texttt{KVCache}, \texttt{CachePolicy}, \texttt{AppendOnlyCache}, \texttt{SlidingWindowCache} & Cache storage and visibility policy. \\
\texttt{Sampler}, \texttt{GreedySampler}, \texttt{TemperatureSampler}, \texttt{TopKSampler}, \texttt{TopPSampler} & Logit-to-token policies. \\
\texttt{GenerationLoop}, \texttt{Prefill}, \texttt{Decode} & Autoregressive scheduling. \\
\texttt{Validator}, \texttt{GoldenVectorExporter}, \texttt{BenchmarkRunner}, \texttt{Profiler} & Correctness, reproducibility, measurement, and diagnostics. \\
\bottomrule
\end{longtable}}


\begin{intuitionbox}{Intuition}
Interfaces should expose invariants, not implementation accidents.
\end{intuitionbox}

\begin{definitionbox}{Mathematical or contract statement}
The engine vocabulary includes config, loading, tensors, devices, tokenizers, modules, caches, samplers, generation, validation, benchmarking, and profiling.
\end{definitionbox}

\begin{implementationbox}{Implementation interpretation}
Use PyTorch for clarity, C++ for portable systems performance, Rust for ownership discipline, and CUDA as a later backend.
\end{implementationbox}

\begin{verificationbox}{How to verify this works}
\begin{itemize}
\item Check shapes and dtypes at the boundary.
\item Compare deterministic outputs against PyTorch golden vectors.
\item Record tolerance, seed, model artifact, and backend.
\item Do not benchmark until validation passes.
\end{itemize}
\end{verificationbox}

\begin{warningbox}{Common failure modes}
\begin{itemize}
\item Letting framework defaults choose dtype or layout.
\item Testing only a batch size of one.
\item Depending on upstream checkpoint names in runtime code.
\item Treating benchmark output as validation.
\end{itemize}
\end{warningbox}

\begin{chapterexercises}
\item State the central invariant in one sentence.\\
\textit{Hint.} Use the conceptual guide vocabulary.
\item Construct a toy example with small shapes.\\
\textit{Hint.} Use $B=2$, $T=3$, $H=8$ where possible.
\item Name a validation test that would catch a plausible bug.\\
\textit{Hint.} Prefer generated golden vectors to handwritten long arrays.
\end{chapterexercises}
